\begin{proof}[Proof of \cref{obj:prop:annotation-threshold}]
Fix the public parameters in the stated domain. The frontier, oracle, and acquisition constants below are finite real numbers depending only on these parameters. In particular,
\[
\gamma>0,\qquad S>0,\qquad 2\gamma+d>0,\qquad \Delta>0.
\]
For an arbitrary auxiliary-size sequence, write
\[
N_n:=n+m_n\qquad(n\in\mathbb N).
\]
All asymptotic arguments may be restricted to the tail \(n\ge2\).

\begin{enumerate}
\item\label{proof:annotation-threshold-risk-rate}
By \cref{obj:thm:sharp-annotation-frontier,obj:thm:supplied-propensity}, choose constants satisfying
\[
0<c_F\le C_F<\infty,\qquad 0<c_O\le C_O<\infty
\]
such that, uniformly for \(n\ge2\) and \(m\ge0\),
\[
c_F r_*(n,m)\le R(n,m)\le C_F r_*(n,m),
\qquad
c_O r_o(n)\le R^e(n,m)\le C_O r_o(n).
\]
Since
\[
r_*(n,m)\ge r_o(n)=n^{-\gamma/(2\gamma+d)}>0,
\]
these bounds make both risks positive and finite. Dividing the upper bound for one risk by the lower bound for the other gives
\[
\begin{aligned}
0\le \frac{R(n,m)}{R^e(n,m)}
&\le \frac{C_F}{c_O}\frac{r_*(n,m)}{r_o(n)},\\
0\le \frac{r_*(n,m)}{r_o(n)}
&\le \frac{C_O}{c_F}\frac{R(n,m)}{R^e(n,m)}.
\end{aligned}
\]
The same calculation, using the frontier bounds at \((n,m)\) and \((n,0)\), gives
\[
\begin{aligned}
0\le \frac{R(n,m)}{R(n,0)}
&\le \frac{C_F}{c_F}\frac{r_*(n,m)}{r_*(n,0)},\\
0\le \frac{r_*(n,m)}{r_*(n,0)}
&\le \frac{C_F}{c_F}\frac{R(n,m)}{R(n,0)}.
\end{aligned}
\]
The first pair transfers eventual boundedness in both directions. The second pair transfers convergence to zero by squeezing between zero and a fixed multiple of the other ratio. Moreover, positivity of \(r_o(n)\) gives
\[
\frac{r_*(n,m_n)}{r_o(n)}=O(1)
\quad\Longleftrightarrow\quad
r_*(n,m_n)=O(r_o(n)).
\]
Consequently,
\[
\begin{aligned}
\frac{R(n,m_n)}{R^e(n,m_n)}=O(1)
&\quad\Longleftrightarrow\quad r_*(n,m_n)=O(r_o(n)),\\
\frac{R(n,m_n)}{R(n,0)}\longrightarrow0
&\quad\Longleftrightarrow\quad
\frac{r_*(n,m_n)}{r_*(n,0)}\longrightarrow0.
\end{aligned}
\]

If the rate comparison holds, there is a real constant \(K_{\mathrm{rate}}\) such that
\[
r_*(n,m_n)\le K_{\mathrm{rate}}r_o(n)
\qquad\text{eventually}.
\]
Set
\[
\zeta:=\max\{K_{\mathrm{rate}},1\}.
\]
Then \(1\le\zeta<\infty\), and the defining inequality in
\cref{obj:def:annotation-region} holds eventually. Conversely, eventual membership in that region supplies the rate comparison with constant \(\zeta\). Thus
\[
r_*(n,m_n)=O(r_o(n))
\quad\Longleftrightarrow\quad
\exists\,\zeta\in[1,\infty)\ 
\text{such that }(n,m_n)\in\mathcal B_*(\zeta)
\text{ eventually}.
\]
% lean: CausalSmith.Stat.TwosamplePointcateAnnotationFrontier.annotation_risk_rate_equivalences

\item
For the extended nonnegative lower limit, the tail characterization is
\[
0<\liminf_{n\to\infty}\frac{N_n}{n^{q_*}}
\quad\Longleftrightarrow\quad
\exists\,c_{\mathrm{count}}>0\
\text{such that }
c_{\mathrm{count}}\le\frac{N_n}{n^{q_*}}
\text{ eventually}.
\]
Indeed, a positive lower limit admits a smaller positive finite real number, which bounds every sufficiently late term from below; this also applies when the lower limit is \(+\infty\). Conversely, an eventual positive lower bound bounds the lower limit from below.

We now translate annotation-region membership into precisely such a bound. Directly from the definitions,
\[
q_*=\frac{S_{\mathrm{crit}}}{S},
\qquad
\Delta=(2\gamma+d)(1+q_*).
\]
For \(n\ge2\), this yields
\[
r_o(n)^{-\Delta/\gamma}
=n^{\Delta/(2\gamma+d)}
=n^{1+q_*}
=n\,n^{q_*}.
\]
Let \(\zeta\ge1\). Since \(r_o(n)\le\zeta r_o(n)\), the maximum defining \(r_*\) and
\cref{obj:def:annotation-region} give
\[
\begin{aligned}
(n,m)\in\mathcal B_*(\zeta)
&\quad\Longleftrightarrow\quad
(nN)^{-\gamma/\Delta}\le\zeta r_o(n)\\
&\quad\Longleftrightarrow\quad
(\zeta r_o(n))^{-\Delta/\gamma}\le nN\\
&\quad\Longleftrightarrow\quad
\zeta^{-\Delta/\gamma}\le\frac{N}{n^{q_*}}.
\end{aligned}
\]
Here the second equivalence uses the strictly decreasing positive-power map with exponent \(-\gamma/\Delta\), whose inverse has exponent \(-\Delta/\gamma\); all bases are positive.

Eventual membership therefore supplies the positive lower bound
\[
c_{\mathrm{count}}:=\zeta^{-\Delta/\gamma}>0.
\]
In the reverse direction, suppose an eventual lower bound \(c_{\mathrm{count}}>0\) is given, and choose
\[
\zeta:=\max\left\{1,c_{\mathrm{count}}^{-\gamma/\Delta}\right\}.
\]
Then
\[
1\le\zeta<\infty,
\qquad
\zeta^{-\Delta/\gamma}
\le
\left(c_{\mathrm{count}}^{-\gamma/\Delta}\right)^{-\Delta/\gamma}
=c_{\mathrm{count}}.
\]
The displayed count characterization gives eventual membership in
\(\mathcal B_*(\zeta)\). Combining this with \cref{proof:annotation-threshold-risk-rate} proves all four oracle-order equivalences.
% lean: CausalSmith.Stat.TwosamplePointcateAnnotationFrontier.annotationRegion_liminf_iff

\item
To characterize strict improvement, first suppose \(S<S_{\mathrm{crit}}\). The identities above imply
\[
q_*=\frac{S_{\mathrm{crit}}}{S}>1,
\qquad
\Delta=(2\gamma+d)(1+q_*)>2(2\gamma+d),
\]
and hence
\[
\frac{2\gamma}{\Delta}<\frac{\gamma}{2\gamma+d}.
\]
Define the positive decay exponent
\[
\eta_{\mathrm{dec}}
:=\frac{\gamma}{2\gamma+d}-\frac{2\gamma}{\Delta}>0.
\]
For \(n\ge2\), the inequality \(q_*>1\) gives
\[
n=n^1\le n^{q_*}.
\]
The lower-count branch in
\cref{obj:thm:sharp-annotation-frontier}, applied with \(m=0\), therefore gives
\[
r_*(n,0)=(n^2)^{-\gamma/\Delta}
=n^{-2\gamma/\Delta}.
\]
Dividing the two terms in the maximum defining \(r_*(n,m_n)\) by this positive rate yields
\[
\frac{r_o(n)}{n^{-2\gamma/\Delta}}
=n^{-\eta_{\mathrm{dec}}},
\qquad
\frac{(nN_n)^{-\gamma/\Delta}}{n^{-2\gamma/\Delta}}
=\left(\frac{N_n}{n}\right)^{-\gamma/\Delta}.
\]
Thus
\[
\frac{r_*(n,m_n)}{r_*(n,0)}
=
\max\left\{
n^{-\eta_{\mathrm{dec}}},
\left(\frac{N_n}{n}\right)^{-\gamma/\Delta}
\right\}.
\]

The second term tends to zero exactly when \(N_n/n\to+\infty\). To see the reverse implication explicitly, let
\[
b_{\mathrm{thr}}\in\mathbb R,
\qquad
x_{\mathrm{thr}}:=\max\{1,b_{\mathrm{thr}}\}>0.
\]
If the second term tends to zero, then eventually
\[
\left(\frac{N_n}{n}\right)^{-\gamma/\Delta}
\le x_{\mathrm{thr}}^{-\gamma/\Delta}.
\]
Inverting this decreasing power gives
\[
\frac{N_n}{n}
\ge
\left(x_{\mathrm{thr}}^{-\gamma/\Delta}\right)^{-\Delta/\gamma}
=x_{\mathrm{thr}}
\ge b_{\mathrm{thr}}.
\]
This proves divergence to \(+\infty\). The forward implication follows from the convergence of a strictly negative power to zero at \(+\infty\). Since
\[
\frac{N_n}{n}=1+\frac{m_n}{n}\qquad(n\ge2),
\]
divergence of the total-to-labeled ratio is equivalent to divergence of the auxiliary-to-labeled ratio: an eventual lower bound at any real threshold for either ratio gives the corresponding bound for the other after shifting the threshold by one.

The first term \(n^{-\eta_{\mathrm{dec}}}\) tends to zero. Both terms are nonnegative, so convergence of their maximum to zero implies convergence of the second term by squeezing; conversely, convergence of both terms implies convergence of their maximum. We have proved, in this regime,
\[
\frac{r_*(n,m_n)}{r_*(n,0)}\longrightarrow0
\quad\Longleftrightarrow\quad
\frac{m_n}{n}\longrightarrow+\infty.
\]

For the complementary case \(S\ge S_{\mathrm{crit}}\),
\cref{obj:thm:sharp-annotation-frontier} gives, for every \(n\ge2\),
\[
r_*(n,m_n)=r_*(n,0)=r_o(n)>0,
\qquad
\frac{r_*(n,m_n)}{r_*(n,0)}=1.
\]
The ratio therefore has limit one, which excludes convergence to zero. Combining the two cases and the risk-to-rate equivalence proves
\[
\frac{R(n,m_n)}{R(n,0)}\longrightarrow0
\quad\Longleftrightarrow\quad
\frac{r_*(n,m_n)}{r_*(n,0)}\longrightarrow0
\quad\Longleftrightarrow\quad
S<S_{\mathrm{crit}}
\ \text{and}\ 
\frac{m_n}{n}\longrightarrow+\infty.
\]
% lean: CausalSmith.Stat.TwosamplePointcateAnnotationFrontier.annotation_rate_improvement_iff

\item
The claimed exponent comparison follows from the same positive-denominator identity:
\[
S<S_{\mathrm{crit}}
\quad\Longrightarrow\quad
q_*=\frac{S_{\mathrm{crit}}}{S}>1,
\qquad S>0.
\]
% lean: CausalSmith.Stat.TwosamplePointcateAnnotationFrontier.qStar_gt_one

\item
Suppose \(m_n/n=O(1)\). Choose a real constant \(K_{\mathrm{aux}}\) satisfying
\[
\frac{m_n}{n}\le K_{\mathrm{aux}}
\qquad\text{eventually},
\]
and define
\[
K_{\mathrm{aux}}^+:=\max\{K_{\mathrm{aux}},0\},
\qquad
M_{\mathrm{aux}}
:=(1+K_{\mathrm{aux}}^+)^{\gamma/\Delta}\ge1.
\]
Eventually \(m_n\le K_{\mathrm{aux}}^+n\), and consequently
\[
n^2\le nN_n\le(1+K_{\mathrm{aux}}^+)n^2.
\]
Taking the negative power \(-\gamma/\Delta\) gives
\[
(nN_n)^{-\gamma/\Delta}\le(n^2)^{-\gamma/\Delta},
\]
and
\[
M_{\mathrm{aux}}^{-1}(n^2)^{-\gamma/\Delta}
\le(nN_n)^{-\gamma/\Delta},
\qquad
(n^2)^{-\gamma/\Delta}
\le M_{\mathrm{aux}}(nN_n)^{-\gamma/\Delta}.
\]
The oracle term is common to both rates and satisfies
\[
r_o(n)\le M_{\mathrm{aux}}r_o(n).
\]
Taking maxima therefore gives
\[
r_*(n,m_n)\le r_*(n,0)
\le M_{\mathrm{aux}}r_*(n,m_n)
\qquad\text{eventually}.
\]
Using the risk sandwiches,
\[
\begin{aligned}
R(n,m_n)
&\le C_F r_*(n,m_n)
\le C_F r_*(n,0)
\le \frac{C_F}{c_F}R(n,0),\\
R(n,0)
&\le C_F r_*(n,0)
\le C_F M_{\mathrm{aux}}r_*(n,m_n)
\le \frac{C_F M_{\mathrm{aux}}}{c_F}R(n,m_n).
\end{aligned}
\]
These eventual inequalities prove both stated big-\(O\) comparisons.
% lean: CausalSmith.Stat.TwosamplePointcateAnnotationFrontier.bounded_auxiliary_risk_order

\item
If \(S\ge S_{\mathrm{crit}}\), the frontier rate equals \(r_o(n)\). For every \(n\ge2\) and \(m\ge0\), the frontier upper bound and the oracle lower bound consequently give
\[
R(n,m)
\le C_F r_o(n)
=\frac{C_F}{c_O}\bigl(c_O r_o(n)\bigr)
\le\frac{C_F}{c_O}R^e(n,m).
\]
Dividing by the positive oracle risk yields
\[
0\le\frac{R(n,m)}{R^e(n,m)}
\le\frac{C_F}{c_O}.
\]
This uniform bound proves the oracle-order comparison for every auxiliary-size sequence.
% lean: CausalSmith.Stat.TwosamplePointcateAnnotationFrontier.smooth_oracle_risk_ratio

\item
For the finite-sample conclusions,
\cref{obj:thm:sharp-annotation-frontier} supplies, simultaneously for all \(n\ge2\) and \(m\ge0\),
\[
c_F r_*(n,m)\le R(n,m)
\le\sup_{P\in\mathcal M}\mathcal R_{n,m}(T_*,P)
\le C_F r_*(n,m).
\]
In particular, its upper bound follows by comparing the minimax value with the worst-case risk of \(T_*\). The same cited statement supplies exactly
\[
r_*(n,m)=
\begin{cases}
r_o(n),&S\ge S_{\mathrm{crit}},\\
(nN)^{-\gamma/\Delta},
&S<S_{\mathrm{crit}},\quad N\le n^{q_*},\\
r_o(n),
&S<S_{\mathrm{crit}},\quad N\ge n^{q_*},
\end{cases}
\]
together with
\[
N=n^{q_*}
\quad\Longrightarrow\quad
(nN)^{-\gamma/\Delta}=r_o(n).
\]
Thus \(c_F,C_F\) furnish the asserted finite-sample constants, including the equality of the two count branches at their boundary.
% lean: CausalSmith.Stat.TwosamplePointcateAnnotationFrontier.sharp_annotation_frontier

\item
For \(n\ge2\), specialize to \(m=0\). If \(S<S_{\mathrm{crit}}\), then \(q_*>1\), so
\[
N=n\le n^{q_*},
\qquad
r_*(n,0)=(n^2)^{-\gamma/\Delta}
=n^{-2\gamma/\Delta}.
\]
If \(S\ge S_{\mathrm{crit}}\), the first branch of
\cref{obj:thm:sharp-annotation-frontier} gives \(r_*(n,0)=r_o(n)\). Hence
\[
r_*(n,0)=
\begin{cases}
n^{-2\gamma/\Delta},&S<S_{\mathrm{crit}},\\
r_o(n),&S\ge S_{\mathrm{crit}}.
\end{cases}
\]
% lean: CausalSmith.Stat.TwosamplePointcateAnnotationFrontier.supervised_rate_branches

\item
Define the average nuisance smoothness by
\[
s_{\mathrm{pub}}:=\frac S2>0.
\]
By \cref{obj:lem:published-cate-benchmark}, algebra with positive denominators gives
\[
\frac{d/4}{1+d/(2\gamma)}
=\frac{S_{\mathrm{crit}}}{2},
\qquad
\frac{1}{1+d/(2\gamma)+d/(4s_{\mathrm{pub}})}
=\frac{2\gamma}{\Delta},
\qquad
\frac{1}{2+d/\gamma}
=\frac{\gamma}{2\gamma+d}.
\]
In particular,
\[
s_{\mathrm{pub}}<\frac{d/4}{1+d/(2\gamma)}
\quad\Longleftrightarrow\quad
S<S_{\mathrm{crit}}.
\]
For \(n\ge2\), these identities identify the numerical benchmark's strict split and its two exponents with the supervised branches just established.

The remaining sample sizes are \(n=0\) and \(n=1\). At \(n=1\), all power bases equal one, so both the benchmark and \(r_*(1,0)\) equal one. At zero, the numerical formulas use the convention
\[
0^{t_{\mathrm{pow}}}=0
\qquad
\bigl(t_{\mathrm{pow}}\in\mathbb R\setminus\{0\}\bigr).
\]
All exponents occurring here are strictly negative:
\[
-\frac{\gamma}{2\gamma+d}<0,
\qquad
-\frac{\gamma}{\Delta}<0,
\qquad
-\frac{2\gamma}{\Delta}<0.
\]
Thus both branches of the benchmark equal zero at \(n=0\), and the two terms in the maximum defining \(r_*(0,0)\) also equal zero. This proves the benchmark equality for every \(n\in\mathbb N\), with the stated formulas for \(n\ge1\).
% lean: CausalSmith.Stat.TwosamplePointcateAnnotationFrontier.published_benchmark_eq_supervised_rate

\item
For acquisition, define the positive exponents
\[
p_0:=\frac{2\gamma+d}{\gamma}
=2+\frac d\gamma,
\qquad
p_1:=\frac{\Delta}{\gamma}
=2+\frac d\gamma+\frac dS.
\]
For positive \(\eta\), the power-scaling identity is
\[
\left(\frac{\eta}{c_{\mathrm{scale}}}\right)^{-p_{\mathrm{scale}}}
=
c_{\mathrm{scale}}^{p_{\mathrm{scale}}}
\eta^{-p_{\mathrm{scale}}}
\qquad
\left(
c_{\mathrm{scale}}\in\{c_F,C_F\},\
p_{\mathrm{scale}}\in\{p_0,p_1\}
\right).
\]
Choose
\[
c_{\mathrm{acq}}
:=\min\{c_F^{p_0},c_F^{p_1}\}>0,
\qquad
C_{\mathrm{acq}}
:=\max\{C_F^{p_0},C_F^{p_1}\}>0.
\]

For any positive \(\eta\), \(n\ge2\), and \(m\ge0\), the defining maximum in
\cref{obj:def:sharp-rate} gives
\[
\begin{aligned}
r_*(n,m)\le\eta
&\quad\Longleftrightarrow\quad
n^{-\gamma/(2\gamma+d)}\le\eta
\ \text{and}\
(nN)^{-\gamma/\Delta}\le\eta\\
&\quad\Longleftrightarrow\quad
\eta^{-p_0}\le n
\ \text{and}\
\eta^{-p_1}\le nN.
\end{aligned}
\]
The second equivalence again inverts strictly decreasing powers of positive bases.

Now let \(0<\eta<1\), and suppose a Borel decision \(T\) satisfies
\[
\mathcal R_{n,m}(T,P)\le\eta
\qquad(P\in\mathcal M).
\]
By \cref{obj:def:minimax} and the frontier lower bound,
\[
c_F r_*(n,m)\le R(n,m)
\le\sup_{P\in\mathcal M}\mathcal R_{n,m}(T,P)
\le\eta.
\]
Therefore \(r_*(n,m)\le\eta/c_F\). Apply the exact rate equivalence at the positive scale \(\eta/c_F\), and then the scaling identity, to obtain
\[
c_F^{p_0}\eta^{-p_0}\le n,
\qquad
c_F^{p_1}\eta^{-p_1}\le nN.
\]
The choice of the minimum gives the common necessity constant:
\[
c_{\mathrm{acq}}\eta^{-p_0}
\le c_F^{p_0}\eta^{-p_0}\le n,
\qquad
c_{\mathrm{acq}}\eta^{-p_1}
\le c_F^{p_1}\eta^{-p_1}\le nN.
\]

Conversely, suppose
\[
C_{\mathrm{acq}}\eta^{-p_0}\le n,
\qquad
C_{\mathrm{acq}}\eta^{-p_1}\le nN.
\]
By the choice of the maximum and the scaling identity,
\[
\left(\frac{\eta}{C_F}\right)^{-p_0}
=C_F^{p_0}\eta^{-p_0}
\le C_{\mathrm{acq}}\eta^{-p_0}\le n,
\]
and
\[
\left(\frac{\eta}{C_F}\right)^{-p_1}
=C_F^{p_1}\eta^{-p_1}
\le C_{\mathrm{acq}}\eta^{-p_1}\le nN.
\]
The exact equivalence at scale \(\eta/C_F\) gives
\[
r_*(n,m)\le\frac{\eta}{C_F}.
\]
By \cref{obj:thm:sharp-annotation-frontier}, the decision \(T_*\) from
\cref{obj:def:sharp-decision} is Borel measurable, and, for every \(P\in\mathcal M\),
\[
\mathcal R_{n,m}(T_*,P)
\le\sup_{P'\in\mathcal M}\mathcal R_{n,m}(T_*,P')
\le C_F r_*(n,m)
\le\eta.
\]
Thus \(c_{\mathrm{acq}},C_{\mathrm{acq}}\) supply the asserted acquisition constants. Finally,
\[
n\le n+m=N
\]
because \(m\ge0\).
% lean: CausalSmith.Stat.TwosamplePointcateAnnotationFrontier.annotation_acquisition
\end{enumerate}
\end{proof}
